\documentclass[12pt]{article}
\usepackage[T1]{fontenc}
\usepackage[utf8]{inputenc}
\usepackage{lmodern}
\usepackage{textcomp}
\usepackage{enumitem}
\usepackage{geometry}
\geometry{margin=1in}
\begin{document}
\begin{center}
Answer Key - Constitutional Law - Graduate Level Assessment 14 \
\small (Answers are ordered exactly as in the question paper.)
\end{center}

\section*{Multiple Choice}
\begin{enumerate}[label=\arabic*.]
\item \textbf{Answer:} A
\\ \textit{Options:} A) Justiciability refers to the limitation of a State to enforce a particular right; B) Justiciability refers to the ethical implications of a particular right; C) Justiciability refers to the just nature or cause of a particular right; D) Justiciability refers to the availability of legal remedies for a particular right; E) Justiciability refers to the obligation of a State to enforce a particular Right
\\ \textit{Note:} Correct option: A - Justiciability refers to the limitation of a State to enforce a particular right
\item \textbf{Answer:} A
\\ \textit{Options:} A) quasi contract; B) unilateral contract; C) secondary party beneficiary contract; D) express contract; E) bilateral contract
\\ \textit{Note:} Correct option: A - quasi contract
\item \textbf{Answer:} A
\\ \textit{Options:} A) 'Difference feminism argues that men and women are fundamentally the same.'; B) Men are unable to comprehend their differences from women.'; C) 'Men and women differ only in physical attributes, not in their thought processes.'; D) 'Men are more inclined towards justice, while women tend to focus on fairness.'; E) There are fundamental differences between individual women.'
\\ \textit{Note:} Correct option: A - 'Difference feminism argues that men and women are fundamentally the same.'
\item \textbf{Answer:} B
\\ \textit{Options:} A) Postmodernism disregards individual stories.; B) Universal values are meaningless.; C) Literature reproduces repression.; D) The concept of truth is objective and undeniable.; E) The law is static.
\\ \textit{Note:} Correct option: B - Universal values are meaningless.
\item \textbf{Answer:} A
\\ \textit{Options:} A) It distinguishes social rules from mere group habits.; B) It defines the judicial function.; C) It illustrates the authority of the legislature.; D) It stresses the relationship between law and justice.
\\ \textit{Note:} Correct option: A - It distinguishes social rules from mere group habits.
\end{enumerate}

\section*{True / False}
\begin{enumerate}[label=\arabic*.]
\setcounter{enumi}{5}
\item \textbf{Answer:} False
\\ \textit{Options:} A) True; B) False
\\ \textit{Note:} LLM-generated false statement. Correct answer: 'Separated law from religion, ethics, sociology and history'.
\item \textbf{Answer:} True
\\ \textit{Options:} A) True; B) False
\\ \textit{Note:} LLM-generated true statement based on MCQ.
\item \textbf{Answer:} False
\\ \textit{Options:} A) True; B) False
\\ \textit{Note:} LLM-generated false statement. Correct answer: 'Warranty of Fitness for Particular Purpose.'.
\item \textbf{Answer:} True
\\ \textit{Options:} A) True; B) False
\\ \textit{Note:} LLM-generated true statement based on MCQ.
\item \textbf{Answer:} False
\\ \textit{Options:} A) True; B) False
\\ \textit{Note:} LLM-generated false statement. Correct answer: 'not hearsay.'.
\end{enumerate}

\section*{Long Form Response}
\begin{enumerate}[label=\arabic*.]
\setcounter{enumi}{10}
\item \textbf{Answer:} A defendant and a victim were driving to work together when the defendant, without provocation, stabbed the victim in the arm. The defendant then pushed him out of the car. The victim fell along the side of the street and fractured his ankle. Unable to move from the roadway, the victim was fatally crushed a half-hour later by a car driven by a drunk driver.. Explain why this is the correct answer and provide supporting evidence. Reference key facts or concepts that support this choice.
\\ \textit{Note:} Long-form derived from MCQ; award credit for stating the correct idea and giving supporting reasoning.
\item \textbf{Answer:} It is an agreement between two or more to commit a crime.. Explain why this is the correct answer and provide supporting evidence. Reference key facts or concepts that support this choice.
\\ \textit{Note:} Long-form derived from MCQ; award credit for stating the correct idea and giving supporting reasoning.
\end{enumerate}

\vfill
\noindent\textit{End of Answer Key}
\end{document}